% TOP_PROBLEM: {"problemId": "problem.corona-problem-for-the-unit-ball-in-several-complex-variables", "canonicalTitle": "Corona Problem for the Unit Ball in Several Complex Variables", "releaseRank": 371, "primaryDomain": "analysis_pde", "primaryDomainLabel": "Analysis & PDE", "displayStatus": "open_with_solved_subcases", "releaseStatus": "open_with_solved_subcases"}
% !TeX program = lualatex
% ATTEMPT_STATUS: unresolved
% Model: GPT-6 Astra Ultra; continuation dated 27 September 2026.
\documentclass[11pt]{article}
\usepackage[margin=25mm]{geometry}
\usepackage{fontspec}
\setmainfont{Latin Modern Roman}
\usepackage{amsmath,amssymb,mathtools}
\usepackage{unicode-math}
\setmathfont{Latin Modern Math}
\usepackage{xurl,hyperref}
\hypersetup{colorlinks=true,urlcolor=blue}
\setcounter{secnumdepth}{0}
\setlength{\parindent}{0pt}
\setlength{\parskip}{5pt}
\setlength{\emergencystretch}{3em}
\begin{document}
\section*{\texorpdfstring{Corona Problem for the Unit Ball in Several Complex Variables}{Corona Problem for the Unit Ball in Several Complex Variables}}\label{371}
\subsection{Definitions and mathematical statement}
Let $B^n=\{z\in{\mathbb{C}}^n:\sum_i|z_i|^2<1\}$ and let $H^\infty(B^n)$ be its algebra of bounded scalar holomorphic functions. For each $n\ge2$ and finite $m\ge1$, assume $f_1,\ldots,f_m\in H^\infty(B^n)$ have a uniform lower bound
\[
 \sum_{j=1}^m|f_j(z)|\ge\delta>0\qquad(z\in B^n).
\]
The question is whether there always exist $g_j\in H^\infty(B^n)$ such that $\sum_j f_jg_j=1$ identically. This is solvability in the bounded holomorphic algebra, not merely in the algebra of all holomorphic functions. The Euclidean ball, rather than the polydisc, is the selected domain. No boundary regularity or optimal estimate on the solution norms is imposed.
\par\smallskip\textit{Formulation and concept references:} \href{https://revista-em.unex.es/index.php/EM/article/download/2411/2144/7535}{[S1]}.

\subsection{Short English statement}
On every higher-dimensional complex unit ball, do bounded holomorphic functions uniformly separated from a common zero always generate the constant function one?
\subsection{Research attempt}
\paragraph{The source's boundary regularity does not disappear.}
The publisher's text of [S1], in its introduction and description
of Theorems 3.4 and 3.6, distinguishes results for $A^k$ and
holomorphic Lipschitz data from the original $H^\infty$ corona
problem on higher-dimensional balls. Its broad abstract language
about solving corona problems refers to those specified algebras.
It is not used here as a verified proof for arbitrary bounded
holomorphic data with no boundary regularity. The publisher's full
PDF fetch was unavailable in this audit; the indexed primary
introduction was accessible, and this scope distinction is what
was checked. No conclusion relies on an uninspected proof in that
paper.

The approach below proves the result for the ball algebra directly,
solves the regularized problems for general data, and identifies
the exact uniform estimate needed to pass to the limit. A separate
two-generator calculation exhibits the bounded $\bar\partial$
problem hidden in the usual correction method.

\paragraph{Two immediate subcases and a stability lemma.}
If $m=1$, the lower bound is $|f_1|\geq\delta$, so
$g_1=1/f_1$ is holomorphic and bounded by $1/\delta$. More
generally, if constants $c_1,\ldots,c_m$ satisfy
\[
 \left|\sum_j c_jf_j(z)\right|\geq\varepsilon>0
                      \quad(z\in B^n),
\]
then
$g_j=c_j/(\sum_i c_if_i)$ is a bounded holomorphic solution.
The corona condition on the vector $f$ does not itself provide
one constant linear combination separated from zero; its direction
can vary across the domain.

Suppose a tuple $f$ already has a bounded solution $g$. If another
tuple $\widetilde f$ satisfies
\[
 q:=\left\|\sum_j(\widetilde f_j-f_j)g_j\right\|_\infty<1,
\]
then
\[
 \widetilde g_j=
 \frac{g_j}{1+\sum_i(\widetilde f_i-f_i)g_i}
 \tag{371.1}
\]
is a solution for $\widetilde f$, with
$\|\widetilde g_j\|_\infty\leq\|g_j\|_\infty/(1-q)$.
The denominator is invertible by its uniformly convergent geometric
series in $H^\infty$. Thus solvable tuples form an open set in
the product norm topology. This is a genuine perturbative result,
but it requires the norm of an existing solution to estimate the
size of that neighborhood.

\paragraph{A full ball-algebra result from its characters.}
Let $A(\overline B^n)$ be the Banach algebra of functions holomorphic
on $B^n$ and continuous on its closure, with the supremum norm.
We first show every character is evaluation at a point of the
closed ball. A character means a nonzero multiplicative linear
functional; in a unital complex Banach algebra it has norm one.
This follows because its value on an element lies in that
element's spectrum, bounded by the element's norm, and its value
on one is one.

If $\chi$ is a character, set $\lambda_j=\chi(z_j)$. For every
$c\in\mathbb C^n$,
\[
 \left|\sum_jc_j\lambda_j\right|
       \leq\left\|\sum_jc_jz_j\right\|_\infty=\|c\|_2.
\]
Choosing $c_j=\overline\lambda_j/\|\lambda\|_2$ when
$\lambda\ne0$ gives $\|\lambda\|_2\leq1$.
Multiplicativity makes $\chi(p)=p(\lambda)$ for every polynomial.

Polynomials are uniformly dense in this ball algebra. Indeed, for
$f\in A(\overline B^n)$, the radial dilates $f_r(z)=f(rz)$
converge uniformly to $f$ by uniform continuity on the closed ball.
For fixed $r<1$, $f_r$ is holomorphic on a larger ball; its Taylor
polynomials at zero converge uniformly on $\overline B^n$.
For example, group the Taylor series by homogeneous degree and
use Cauchy estimates on an intermediate radius larger than one
to obtain a geometric majorant. Thus continuity of $\chi$ gives
$\chi(f)=f(\lambda)$ for every ball-algebra function.

Now suppose the $f_j$ belong to this algebra and satisfy the
printed lower bound. Continuity extends it to the closed ball.
If their generated ideal were proper, it would lie in a maximal
ideal. A maximal ideal of a commutative unital Banach algebra is
closed: if it were dense, it would contain an element within
distance one of the unit, hence an invertible element, contradicting
properness. The quotient is a complex Banach division algebra
and therefore is the scalars, giving a character annihilating
all $f_j$. The character just identified is evaluation at some
$\lambda\in\overline B^n$, contrary to the lower bound.
Hence the ideal is all of $A(\overline B^n)$, and
\[
 \sum_jf_jg_j=1\quad\hbox{with }g_j\in A(\overline B^n).
 \tag{371.2}
\]
This proves a complete several-variable subcase. Its key input,
uniform approximation by polynomials, used boundary continuity.

\paragraph{Every radial regularization is solvable.}
For arbitrary bounded holomorphic corona data, put
$f_{j,r}(z)=f_j(rz)$ for $0<r<1$. These functions extend
holomorphically to a neighborhood of the closed ball and satisfy
the same lower bound $\sum_j|f_{j,r}|\geq\delta$ there on
$\overline B^n$. The preceding ball-algebra proof therefore gives
bounded holomorphic solutions $g_{j,r}$ of
\[
 \sum_j f_j(rz)g_{j,r}(z)=1.
 \tag{371.3}
\]
This does not yet solve the original tuple. The proof gives no
bound on the solution norms uniform as $r\uparrow1$.

If such a bound were available along some sequence $r_\ell\uparrow1$,
Montel's theorem and a diagonal subsequence would give locally
uniform convergence of every $g_{j,r_\ell}$ to a holomorphic
$g_j$. The common uniform bound would persist in the limit.
Since $f_j(r_\ell z)\to f_j(z)$ locally uniformly, (371.3)
would pass to the limit and solve the requested problem. Thus
we have the precise sufficient estimate
\[
 \sup_\ell\max_j\|g_{j,r_\ell}\|_\infty<\infty.
 \tag{371.4}
\]
There is no requirement here for an optimal bound in $\delta$;
even a finite bound depending on the original tuple would suffice.
The obstacle is uniformity in the approach to the boundary.

\paragraph{Why norm approximation cannot replace that estimate.}
In general $f(r\,\cdot)$ does not converge to $f$ in the
$H^\infty$ norm. For an explicit example, take
\[
 f(z)=\exp\left(-\frac{1+z_1}{1-z_1}\right).
 \tag{371.5}
\]
The real part of $(1+z_1)/(1-z_1)$ is positive on the unit
disc, so $|f|\leq1$ on the ball. Along real $z_1\uparrow1$
it tends to zero. On the other hand, choose
$w=1+it$ with $t\to\infty$ and set
$z_1=(w-1)/(w+1)$. Then $|z_1|<1$, $z_1\to1$, and
$f=e^{-1}e^{-it}$ has no limit. Thus $f$ has no continuous
extension at the boundary point $(1,0,\ldots,0)$.

Every fixed radial dilate belongs to the ball algebra. Were those
dilates to converge in supremum norm, their limit would also
extend continuously to the closed ball, since that algebra is
closed in the supremum norm. The example contradicts this.
Consequently the open-neighborhood argument (371.1) cannot be
applied using local uniform convergence alone. Moreover its
allowed perturbation size depends on $\|g_{j,r}\|_\infty$,
which may deteriorate at exactly the same boundary limit.

\paragraph{The bounded correction problem for two generators.}
For clarity take $m=2$ and write
$S=|f_1|^2+|f_2|^2$. The corona assumption gives
$S\geq\delta^2/2$. The smooth, generally nonholomorphic functions
\[
 h_j=\bar f_j/S
\]
satisfy $f_1h_1+f_2h_2=1$ and are uniformly bounded. Define the
$(0,1)$-form
\[
 \alpha=h_1\bar\partial h_2-h_2\bar\partial h_1.
 \tag{371.6}
\]
Differentiating the scalar identity gives
$f_1\bar\partial h_1+f_2\bar\partial h_2=0$, hence
\[
 f_1\alpha=\bar\partial h_2,
       \qquad f_2\alpha=-\bar\partial h_1.
 \tag{371.7}
\]
Since $f_1,f_2$ never vanish simultaneously, the same relation
implies $\bar\partial h_1\wedge\bar\partial h_2=0$ pointwise.
It follows by differentiating (371.6) that
$\bar\partial\alpha=0$.

If one could solve $\bar\partial v=\alpha$ with $v$ bounded,
then
\[
 g_1=h_1+f_2v,\qquad g_2=h_2-f_1v
 \tag{371.8}
\]
would be bounded and $\bar\partial$-closed by (371.7), hence
holomorphic. Their Bezout identity holds because the correction
terms cancel algebraically. This is an exact reduction, including
the signs of both corrections.

The form (371.6) involves derivatives of bounded holomorphic data.
Cauchy estimates only give derivative bounds deteriorating toward
the boundary, typically with an inverse distance factor. The
lower bound on $S$ controls denominators but does not remove
those derivatives. A general bounded $\bar\partial$ solution
estimate for this structured family is precisely the difficult
analytic input; it is not obtained just from
$\bar\partial\alpha=0$.

An $L^2$ solution would not suffice. The holomorphic function
$(1-z_1)^{-1/2}$, with its analytic branch in the ball, is
unbounded but lies in $L^2(B^n)$. Indeed slicing by $z_1=w$
reduces its squared norm to an integral bounded by a constant
times $\int_{|w|<1}|1-w|^{-1}\,dA(w)$, which is finite by
polar integration around one. Thus square integrability does
not upgrade to boundedness even for a holomorphic correction.

\paragraph{An equivalent maximal-ideal obstruction.}
For the full Banach algebra $H^\infty(B^n)$, let its character
space have the topology of pointwise convergence on the algebra.
The corona assertion for all finite tuples is equivalent to
density of the interior evaluation characters in that space.

If evaluations are dense, a tuple with the uniform pointwise
lower bound has the same lower bound at every character by
continuity. No character can annihilate it, and the maximal-ideal
argument gives a Bezout solution. Conversely, if a character
$\chi$ is outside the closure of evaluations, some basic
neighborhood of it excludes every evaluation. Such a neighborhood
uses finitely many functions $F_1,\ldots,F_m$ and a positive
tolerance. The tuple $F_j-\chi(F_j)$ then has a uniform positive
sum of absolute values at every interior point, while all its
entries vanish at $\chi$. It cannot generate one.

The coordinate-character argument for the ball algebra does not
identify all characters of $H^\infty$: polynomial density in
the required norm is missing, as (371.5) demonstrates. Assuming
every character is evaluation, or even assuming evaluation
density without proof, would insert the desired conclusion at
the critical step.

\paragraph{Outcome and verification scope.}
The script samples the correction identities (371.7)--(371.8)
and inverse perturbation formula numerically at three complex points
for one test tuple. Their algebraic proofs are given above.
These samples do not
provide a bounded solution operator for (371.6). The proved
subcases are one generator, a uniformly invertible constant
combination, small perturbations of a solved tuple, and all
boundary-continuous ball-algebra data. Every radial regularization
of a general tuple is also solved.

The first unresolved step is the uniform bound (371.4), or an
equivalent bounded correction estimate for general corona data.
The source's additional boundary-regularity hypotheses do not
supply it for arbitrary $H^\infty$ functions. The selected
unrestricted ball problem remains unresolved by this attempt.
\subsection{Sources}
\begin{itemize}
\item[S1]S. R. Patel, Extracta Mathematicae 41 (2026), 159–184.
\url{https://revista-em.unex.es/index.php/EM/article/download/2411/2144/7535}.
\textit{Formulation source.}
\end{itemize}
\end{document}
